\section {Thirteen coins (Information theory)}
% follows: twelve_coins

The same as the previous riddle, but with $13$ coins: one of them is fake and is either heavier or lighter than the others.
With $3$ weighings, can you find the fake coin and tell whether it is heavier or lighter?
If not, what can you do?

\subsection{Solution}

Counting alone does not rule it out: there are $26$ possibilities and $27$ results of three weighings. But it cannot be done.

Suppose the first weighing puts $k$ coins on each side.
If the scale balances, the fake coin is one of the other $13-2k$ coins, and there are $2(13-2k)$ possibilities left for only $9$ results of the two remaining weighings,
so $26-4k \le 9$, which means $k \ge 5$.
If the left side is heavier, the fake is one of the $k$ coins on the left and heavy, or one of the $k$ coins on the right and light.
That is $2k$ possibilities for $9$ results, so $k \le 4$.
Both cases can happen, so we need both $k \ge 5$ and $k \le 4$, which is impossible.

What can be done is to always find the fake coin, and almost always tell whether it is heavier or lighter.
Weigh $1,2,3,4$ against $5,6,7,8$.
If the scale does not balance, coins $9$ to $13$ are genuine and we go on exactly as in the twelve coin riddle, finding the coin and its weight.
If it balances, the fake coin is one of $9$ to $13$ and coins $1$ to $8$ are genuine. Weigh $9,10,11$ against $1,2,3$.
If this does not balance we know whether the fake is heavy or light, and weighing $9$ against $10$ finds it.
If it balances, the fake is $12$ or $13$. Weigh $12$ against $1$: if it does not balance then $12$ is the fake and we see whether it is heavy or light,
and if it balances then $13$ is the fake, but we never weighed it, so we do not know whether it is heavier or lighter.

If we are also given one extra coin that is known to be genuine, then $13$ coins can be solved completely, since the genuine coin lets the first weighing
put $5$ unknown coins against $4$ unknown coins and the genuine coin.
